% !PRINTED 576
\weberpage{576}
%<<CONT
and consequently, by interchange of $\xi$ with $\xi'$:
\weq{11}{\theta_\xi(\tau_{\xi'}) = \theta_{\xi'}(\tau_\xi),}
and hence, according to Bd. I, \S\ 169:

\begin{quote}
7. \gesperrt{The division field is relatively Abelian with respect to the class field.}
\end{quote}

The numbers $\xi$, taken with respect to the modulus $\mathfrak{m}$, form a group under multiplication, and the relative group of the division field is isomorphic with this group (or at least with a divisor of this group. The irreducibility of $T_{\mathfrak m}$ will be proved further on, whereby the group of $\mathfrak{T}_{\mathfrak m}$ itself is then established).

\begin{center}
\S\ 155. \textbf{Multiplication of the elliptic functions for an}\\
\textbf{odd multiplier.}
\end{center}

For the demonstration of the existence of the division field, the use of the \gesperrt{Weierstrass} function and of the $\tau$-function derived from it is very expedient, and it would be most convenient if the further investigation of this field could also be based upon it. But for that the arithmetic nature of the multiplication formulae is not yet sufficiently known, and it is therefore for the present still necessary to consider the complex multiplication and division of the \gesperrt{Jacobi} elliptic functions as well.

Here we wish to give the multiplication formulae which we derived in \S\ 57 yet a somewhat different shape.

We consider throughout, to begin with, the multiplier $m$ as an \gesperrt{odd} number. We set, if
\weq{1}{\varkappa = \frac{\vartheta_{10}}{\vartheta_{00}}}
denotes the modulus of the elliptic functions, following \gesperrt{Kronecker}\textsuperscript{1)}\begingroup\renewcommand{\thefootnote}{1)}\footnotetext[1]{Zur Theorie der elliptischen Funktionen. Sitzungsbericht der Berliner Akademie vom 29. Juli 1886.}\endgroup:
\weq{2}{\lambda = 4\left(\varkappa + \frac{1}{\varkappa}\right) = 4\,\frac{f(\omega)^4}{f_2(\omega)^4} + 4\,\frac{f_2(\omega)^4}{f(\omega)^4} \quad [\S\ 54, (3)],}
and develop this expression in ascending powers of
\[
q = e^{\pi i \omega}.
\]
